% ============================================================
%  Experiment 8 --- Temperature Control
% ============================================================
\experimentheader{8}{Temperature Control}{Instrumentation --- Temperature Measurement and Control}{Instrumentation Lab}

\begin{objectives}
  \begin{enumerate}
    \item To design and build a signal-conditioning circuit for an NTC thermistor.
    \item To interface the conditioned temperature signal to the PC via the PCIE-1750-U and ADAM-3937.
    \item To implement a computer-based system that displays the measured temperature and performs ON/OFF control of an indicator LED on the trainer.
  \end{enumerate}
\end{objectives}

\begin{equipment}
  \begin{itemize}
    \item NTC thermistor and trainer board with indicator LED
    \item Components for the student-designed signal-conditioning circuit (op-amps, resistors, capacitors)
    \item ADAM-3937 DB37 wiring terminal board
    \item PCIE-1750-U isolated digital I/O card, installed in the host PC
    \item PC with a C/C++ or Python development environment
    \item Digital multimeter
  \end{itemize}
\end{equipment}

\begin{introduction}
  In many industrial processes, temperature must be regulated automatically. Two general strategies are used: \textit{logic (ON/OFF) control}, in which a heater, cooler, or indicator is simply switched on or off when a threshold is crossed, and \textit{fuzzy/proportional control}, which requires a more elaborate definition of ``hot'' or ``cold'' before acting. This experiment uses logic control.
  \vspace*{1.5em}

  An NTC (negative temperature coefficient) thermistor is a resistor whose resistance decreases as its temperature increases. Given the thermistor's resistance--temperature characteristic (from its datasheet), a measured resistance can be converted to a temperature, and vice versa.
  \vspace*{1.5em}

  The signal path for this experiment is: the thermistor's resistance is converted to a voltage by a signal-conditioning circuit (which the student must design), this voltage is digitised and read by the PC via the PCIE-1750-U (through the ADAM-3937 wiring terminal), and software on the PC computes the corresponding temperature, displays it, and drives the indicator LED on the trainer according to a threshold comparison (Figure~\ref{fig:exp8-blockdiagram}).
\end{introduction}

\begin{boxedfigure}
  \centering
  \resizebox{0.98\linewidth}{!}{%
    \begin{tikzpicture}[node distance=10mm and 12mm]
      \node[block] (therm) {NTC\\Thermistor};
      \node[block, right=of therm, minimum width=2.6cm] (cond) {Signal\\Conditioning\\Circuit};
      \node[block, right=of cond] (adamdi) {ADAM-3937\\ (DI)};
      \node[block, right=of adam] (pcie) {PCIE-1750-U};
      \node[block, right=of pcie, minimum width=2.6cm] (pc) {PC};
      \node[block, below=14mm of pc] (adamdo) {ADAM-3937\\ (DO)};
      \node[block, left=of adamdo] (reljt) {Relay \& \\BJT};
      \node[block, below=14mm of adamdi] (led) {Indicator LED\\(trainer)};

      \draw[arr] (therm) -- (cond);
      \draw[arr] (cond) -- (adamdi);
      \draw[arr] (adamdi) -- (pcie);
      \draw[arr] (pcie) -- (pc);
      \draw[arr] (pc) -- (adamdo);
      \draw[arr] (adamdo) -- (reljt);
      \draw[arr] (reljt) -- (led);
    \end{tikzpicture}
  }
  \captionof{figure}{Block diagram of the temperature measurement and control system.}
  \label{fig:exp8-blockdiagram}
\end{boxedfigure}

\begin{procedure}

  \textbf{Activity 1 --- Signal-conditioning circuit design}\vspace{2pt}

  \textit{Objective: design and build a circuit that converts the thermistor's resistance to a voltage suitable for the PCIE-1750-U input.}

  \begin{enumerate}
    \item Obtain the resistance--temperature characteristic for the NTC thermistor from its datasheet, and determine the resistance range corresponding to the temperature range you intend to measure and control.
    \item Design a signal-conditioning circuitthat converts the thermistor's resistance to a voltage. Calculate your component values and justify your design choices, including the expected output voltage range over your chosen temperature range.
    \item Include a voltage-follower (buffer) stage to eliminate loading effects between the conditioning circuit and the downstream measurement hardware.
    \item Build the circuit and verify its output voltage against a reference temperature measurement at two or more points, and check linearity/consistency with your design calculations.
  \end{enumerate}
  \vspace*{1.5em}

  \textbf{Activity 2 --- Interfacing to the DAQ hardware}\vspace{2pt}

  \textit{Objective: connect the conditioned temperature signal and the LED control line to the PC via the ADAM-3937 and PCIE-1750-U.}

  \begin{enumerate}
    \item Connect the output of your signal-conditioning circuit to an input channel on the ADAM-3937 terminal board.
    \item Connect the trainer's indicator LED control line to an isolated digital output channel on the ADAM-3937 via BJT/Relay circuit.
    \item Connect the ADAM-3937 to the PCIE-1750-U card installed in the host PC via the supplied cable.
    \item Verify both the input and output wiring using the DAQ card's bundled configuration/monitoring utility before writing any code: confirm the conditioned voltage reads correctly, and confirm the LED can be toggled from the utility.
  \end{enumerate}
  \vspace*{1.5em}

  \textbf{Activity 3 --- Temperature monitoring and LED control software}\vspace{2pt}

  \textit{Objective: implement software that displays the measured temperature and controls the LED according to a threshold.}

  \begin{enumerate}
    \item In C, C++, or Python, write a program that:
          \begin{itemize}
            \item Periodically reads the conditioned thermistor voltage,
            \item Converts the reading to a temperature using your circuit's calibration from Activity 1,
            \item Displays the current temperature in a simple GUI, updating in real time, and
            \item compares the measured temperature against the threshold and sets the LED outputs accordingly.
          \end{itemize}
    \item Report and justify your design choice for signal conditioning circuit.
    \item Use the flowchart in Figure~\ref{fig:exp8-flowchart} as a starting point for your control logic, adapting the ON/OFF branch to match your chosen design.
  \end{enumerate}

  \begin{boxedfigure}
    \centering
    \resizebox{0.5\linewidth}{!}{%
      \begin{tikzpicture}[node distance=7mm and 10mm]
        \node[startstop] (start) {Start};
        \node[block, below=of start] (init) {Initialise DAQ card};
        \node[block, below=of init] (setpoint) {Read user-specified threshold temperature};
        \node[block, below=of setpoint] (readv) {Acquire thermistor voltage};
        \node[block, below=of readv] (calc) {Convert to temperature};
        \node[decision, below=of calc, minimum width=3.2cm] (dec) {Measured vs.\\threshold?};
        \node[block, below left=12mm and -6mm of dec] (ledon) {Set LED per\\design (case A)};
        \node[block, below right=12mm and -6mm of dec] (ledoff) {Set LED per\\design (case B)};
        \node[block, below=26mm of dec] (disp) {Display temperature \& LED state};

        \draw[arr] (start) -- (init);
        \draw[arr] (init) -- (setpoint);
        \draw[arr] (setpoint) -- (readv);
        \draw[arr] (readv) -- (calc);
        \draw[arr] (calc) -- (dec);
        \draw[arr] (dec.west) -| (ledon) node[pos=0, above, font=\small, xshift=-15pt] {above};
        \draw[arr] (dec.east) -| (ledoff) node[pos=0, above, font=\small, xshift=15pt] {below};
        \draw[arr] (ledon.south) |- ($(disp.north)+(-0.9,0)$);
        \draw[arr] (ledoff.south) |- ($(disp.north)+(0.9,0)$);
        % Loop back: drop below disp, go left past the whole diagram, then up into readv.
        \coordinate (loopY) at ($(disp.south)+(0,-0.6)$);
        \coordinate (loopX) at ($(readv.west)+(-1.8,0)$);
        \draw[arr] (disp.south) -- (disp.south |- loopY)
        -- (loopX |- loopY)
        -- (loopX |- readv.west)
        -- (readv.west);
      \end{tikzpicture}
    }
    \captionof{figure}{Temperature monitoring and LED control loop (ON/OFF direction per student design).}
    \label{fig:exp8-flowchart}
  \end{boxedfigure}

  \begin{enumerate}[resume]
    \item Set the controlled (threshold) temperature to a chosen value and record the measured temperature and LED state at 1-minute intervals for 35 minutes (Table~\ref{tab:exp8-results}). Plot temperature against time.
    \item Repeat for at least two further threshold temperatures. Plot your results.
    \item Discuss your findings, including the responsiveness of the LED control and the accuracy of your temperature conversion.
  \end{enumerate}

  \needspace{4cm}
  \begin{boxedtable}
    \captionof{table}{Temperature and LED-state log.}
    \label{tab:exp8-results}
    \centering
    \begin{tabular}{cccc}
      \toprule
      Threshold temp. & Time (min) & Measured temp. & LED state \\
      \midrule
                      & 1          &                &           \\
                      & 2          &                &           \\
                      & \ldots     &                &           \\
                      & 35         &                &           \\
      \bottomrule
    \end{tabular}
  \end{boxedtable}

\end{procedure}

\begin{reportq}
  \begin{enumerate}
    \item Present your signal-conditioning circuit design (topology, calculations, component values) and its calibration against a reference temperature measurement.
    \item State and justify your chosen LED ON/OFF logic (Activity 3, step 2).
    \item Present and discuss your temperature-vs-time plots for each threshold tested.
    \item What sources of error or lag affect the accuracy and responsiveness of this system (thermistor thermal mass, signal-conditioning bandwidth, software polling rate, etc.)?
  \end{enumerate}
\end{reportq}
